\documentclass[a4paper,11pt]{article}
\usepackage[utf8]{inputenc}
\usepackage[T1]{fontenc}
\usepackage[french]{babel}
\usepackage{graphicx}
\usepackage{amsmath}
\usepackage{geometry}
\usepackage{fancyhdr}
\usepackage{hyperref}
\usepackage[usenames,dvipsnames,svgnames,table]{xcolor}
\usepackage{enumerate}
\usepackage{booktabs}
\usepackage{tcolorbox}
\tcbuselibrary{listings,skins,breakable}
\usepackage{sourcecodepro}
\usepackage{FiraSans}
\usepackage{environ}
\usepackage{tikz}
\usepackage{pgfplots}
\pgfplotsset{compat=1.18}
\usetikzlibrary{arrows.meta,positioning,calc}

\geometry{a4paper, scale=0.8}
\urlstyle{same}

%% ─── Colors ──────────────────────────────────────────────────────
\definecolor{modernblue}{RGB}{66,133,244}
\definecolor{moderngray}{RGB}{95,99,104}
\definecolor{darkgray}{RGB}{60,64,67}
\definecolor{accentorange}{RGB}{255,111,0}
\definecolor{codegreen}{RGB}{15,157,88}
\definecolor{codeblue}{RGB}{25,118,210}
\definecolor{codebg}{RGB}{250,250,250}
\definecolor{codepurple}{RGB}{156,39,176}
\definecolor{niceblue}{RGB}{0,128,172}

%% ─── Code listing: PYTHON ───────────────────────────────────────
\newtcblisting{pythoncode}[1][]{
    listing engine=listings, colback=codebg, colframe=moderngray!50,
    listing only, left=5mm, enhanced, sharp corners=all, boxrule=0.5pt, top=0pt,
    listing options={
        language=Python, basicstyle=\ttfamily\scriptsize,
        keywordstyle=\color{codeblue}\bfseries,
        commentstyle=\color{codegreen}\itshape,
        stringstyle=\color{accentorange},
        numberstyle=\tiny\color{moderngray}, numbers=left, numbersep=8pt,
        breaklines=true, tabsize=4, showstringspaces=false,
        keepspaces=true, columns=flexible,
        literate=
            {é}{{\'{e}}}1 {è}{{\`{e}}}1 {ê}{{\^{e}}}1 {ë}{{\"{e}}}1
            {û}{{\^{u}}}1 {ù}{{\`{u}}}1 {î}{{\^{i}}}1 {ï}{{\"{i}}}1
            {ô}{{\^{o}}}1 {â}{{\^{a}}}1 {à}{{\`{a}}}1 {ä}{{\"{a}}}1
            {ç}{{\c{c}}}1 {Ç}{{\c{C}}}1
            {É}{{\'{E}}}1 {È}{{\`{E}}}1 {Ê}{{\^{E}}}1,
        #1
    }, breakable
}

%% ─── Lesson box ──────────────────────────────────────────────────
\newtcolorbox{lesson}[1][]{
    enhanced, breakable,
    colback=modernblue!5, colframe=modernblue!80,
    colbacktitle=modernblue, coltitle=white,
    fonttitle=\sffamily\bfseries, fontupper=\sffamily\small,
    boxrule=0.5pt, arc=3pt, left=8pt, right=8pt, top=8pt, bottom=8pt,
    attach boxed title to top left={yshift=-2mm, xshift=5mm},
    boxed title style={sharp corners, boxrule=0.5pt},
    title={#1}, before skip=1em, after skip=2em
}

%% ─── Solution toggle ─────────────────────────────────────────────
\newif\ifshowsolutions
\showsolutionsfalse

\ifshowsolutions
    \newtcolorbox{solution}{colback=white, colframe=white, breakable,
        fontupper=\color{red}}
\else
    \NewEnviron{solution}{}{}
\fi

%% ─── Question command ─────────────────────────────────────────────
\newcounter{question}[section]
\newcommand{\q}[1]{%
    \vspace{0.5em}\noindent
    {\color{black}\bfseries\stepcounter{question}\Alph{section}\arabic{question}.}
    \hspace{0.5em}#1\par\vspace{0.3em}
}
\renewcommand{\thesection}{{\Alph{section}}}

%% ─── Header / footer ─────────────────────────────────────────────
\pagestyle{fancy}
\fancyhf{}
\fancyhead[L]{2025-2026}
\fancyhead[C]{\textbf{\courseName}}
\fancyhead[R]{\large GEII -- BUT2}
\fancyfoot[C]{\thepage}
\fancyfoot[L]{\includegraphics[height=3em]{../../../common/logos_couleurs/IUTA.png}}
% \fancyfoot[R]{\href{mailto:prenom.nom@univ-smb.fr}{prenom.nom@univ-smb.fr}}

%% ─── Course metadata ─────────────────────────────────────────────
\newcommand{\courseName}{Outils Mathématiques et Logiciel --- OML3}
\newcommand{\chapterTitle}{TP 2 --- Série de Fourier, forme complexe}

\makeatletter
\renewcommand\maketitle{{\raggedright\begin{center}
    {\Huge\bfseries\@title}
\end{center}}}
\makeatother

\usepackage{amsmath,amsfonts,bm}

%% ─── Operators ───────────────────────────────────────────────────
\DeclareMathOperator*{\argmax}{arg\,max}
\DeclareMathOperator*{\argmin}{arg\,min}
\DeclareMathOperator{\MSE}{MSE}
\DeclareMathOperator{\sign}{sign}
\DeclareMathOperator{\sinc}{sinc}
\DeclareMathOperator{\tr}{tr}

%% ─── Number sets ─────────────────────────────────────────────────
\def\R{{\mathbb{R}}}
\def\C{{\mathbb{C}}}
\def\N{{\mathbb{N}}}
\def\Z{{\mathbb{Z}}}

%% ─── Infinitesimal and units ────────────────────────────────────
\def\d{{\rm d}}
\def\Hz{{\mbox{Hz}}}
\def\m{{\mbox{m}}}

%% ─── Complex / signal ───────────────────────────────────────────
\DeclareMathOperator{\complexj}{\mathsf{j}}   % imaginary unit

%% ─── Probability ─────────────────────────────────────────────────
\newcommand{\Esp}[2]{\mathbb{E}_{#1}\!\left[#2\right]}
\renewcommand{\Pr}{\mathbb{P}}

%% ─── Linear algebra ──────────────────────────────────────────────
\newcommand{\abs}[1]{\left|{#1}\right|}
\newcommand{\norm}[1]{\left\|#1\right\|}
\newcommand{\inner}[2]{\left\langle#1,\,#2\right\rangle}

%% ─── Logic ───────────────────────────────────────────────────────
\newcommand{\then}{\Rightarrow}
\newcommand{\iif}{\Leftrightarrow}

%% ─── Draft annotation ────────────────────────────────────────────
\newcommand{\red}[1]{{\color{red}#1}}

%% ─── Shared figure color (used by common/figures/*/*.tex bodies) ──
\definecolor{niceblue}{RGB}{0,128,172}

%% ─── Rappel header (styled subtitle for reference-card asides) ───
%% Requires \definecolor{niceblue}{...} to be defined by the including document.
\newcommand{\rappelhead}[1]{%
    \bigskip
    {\color{niceblue}\bfseries #1}\\[-0.4em]
    {\color{niceblue!35}\rule{\linewidth}{0.6pt}}
    \smallskip

}


\title{\chapterTitle}
\author{}
\date{}

%% ─────────────────────────────────────────────────────────────────
\begin{document}
\fontfamily{bch}\selectfont
\fontsize{12pt}{14pt}\selectfont

\maketitle
\thispagestyle{fancy}


%% ─── AVANT DE COMMENCER ──────────────────────────────────────────
\section*{Avant de commencer}

Ce TP est écrit en Python. Les fonctions dont vous aurez besoin sont déjà
écrites. Votre travail consiste à exécuter les cellules en modifiant ou
non les paramètres, et à lire les tracés produits.

\textbf{Ce document est le sujet de référence.} Le notebook
\texttt{TP2.ipynb} est un compagnon d'exécution : il ne reprend que
l'essentiel de chaque question, repérée par son numéro (A1, A2, B1,
\dots) correspondant à celui utilisé ici. Gardez ce PDF ouvert pendant
tout le TP.

\textbf{Vos réponses vont sur feuille.} Contrairement au TP1, ce
notebook ne contient pas d'emplacement de réponse : les calculs et
justifications demandés se font sur papier, à rendre à la fin de la
séance.

\subsection*{Configuration de l'environnement}
\begin{enumerate}
    \item Ouvrir \texttt{TP2.ipynb} dans Jupyter.
    \item Exécuter la première cellule (import des fonctions fournies,
          dans \texttt{tp2\_helpers.py}) : rien ne devrait s'afficher,
          c'est normal.
    \item Exécuter ensuite les cellules dans l'ordre, en vous référant aux
          questions ci-dessous.
\end{enumerate}

\vspace{2em}\hrule\vspace{2em}
\newpage


%% ─── SECTION A ──────────────────────────────────────────────────
\section{Forme complexe}

\subsection*{Contexte}
Les séries de Fourier du TP1 s'écrivaient avec des sinus et cosinus.
Cette section fait le lien avec l'écriture complexe
$f(t) = \sum_{n=-\infty}^{+\infty} c_n\, e^{jn\omega_0 t}$, et vérifie sur
un exemple que les deux écritures portent la même information.

\q{Exécuter la cellule de rappel sur les nombres complexes en Python.
Vérifier qu'un réel s'écrit bien comme un complexe de partie imaginaire
nulle (\texttt{2 + 0j}). Calculer $e^{j\theta}$ pour $\theta=2k\pi$,
$(2k+1)\pi$, $\pi/2$ et $\pi/3$ : les valeurs obtenues correspondent-elles
à ce que vous attendiez ?}
\begin{solution}
    $e^{j2k\pi}=1$, $e^{j(2k+1)\pi}=-1$, $e^{j\pi/2}=j$,
    $e^{j\pi/3}=\cos(\pi/3)+j\sin(\pi/3)=\tfrac12+j\tfrac{\sqrt3}{2}$ :
    exactement les valeurs usuelles du cercle trigonométrique, retrouvées
    ici sous forme complexe.
\end{solution}

\q{\textbf{Sur papier.} Le signal carré $k(t)$ a pour coefficients
réels $a_n=0$ et $b_n=4/(n\pi)$ si $n$ impair, $0$ sinon. En utilisant
$c_0=a_0$, $c_n=(a_n-jb_n)/2$ pour $n>0$, $c_{-n}=\overline{c_n}$,
exprimer $c_n$ pour tout $n\in\mathbb Z$.}
\begin{solution}
    $c_0=0$. Pour $n>0$ impair : $c_n = -j\,2/(n\pi)$, et $c_n=0$ si $n$
    pair. Pour $n<0$ : $c_n=\overline{c_{-n}}=+j\,2/(|n|\pi)$ si $n$
    impair (négatif), $0$ sinon.
\end{solution}

\q{Exécuter la cellule qui calcule numériquement les $c_n$ du carré
(\texttt{to\_complex\_coeffs}). Comparer à votre calcul précédent. Les
$c_n$ sont-ils réels, imaginaires purs, ou ni l'un ni l'autre ? Pourquoi ?}
\begin{solution}
    Les valeurs numériques concordent avec A2. Les $c_n$ sont
    imaginaires purs pour tout $n\neq0$ : c'est une conséquence directe
    de l'imparité de $k$ -- $a_n=0$ équivaut exactement à
    $\mathrm{Re}(c_n)=0$ pour $n\neq0$.
\end{solution}

\begin{lesson}[Fonction \texttt{plot\_spectrum\_complex}]
    \texttt{plot\_spectrum\_complex(c, freq=False, T=None)} trace
    $|c_n|$ en fonction du rang $n$ (spectre \emph{bilatéral}, $n$ positif
    et négatif). Avec \texttt{freq=True} et \texttt{T} fourni, l'axe des
    abscisses devient la fréquence $n/T$ au lieu du rang.
\end{lesson}

\q{Tracer le spectre bilatéral du carré. Comment décroît $|c_n|$ avec
$n$ ?}
\begin{solution}
    $|c_n|=|b_n|/2$ pour $n\neq0$ : même décroissance en $1/n$ qu'au TP1,
    à un facteur $2$ près (l'énergie de chaque raie du spectre unilatéral
    du TP1 est répartie sur deux raies, $+n$ et $-n$, du spectre
    bilatéral).
\end{solution}

\q{Retracer le spectre avec \texttt{freq=True}. Qu'est-ce qui change sur
l'axe des abscisses ? Le spectre bilatéral (en $n$, ou en $n/T$) contient-il
plus d'information que le spectre unilatéral du TP1, ou seulement la
même information présentée autrement ?}
\begin{solution}
    L'axe passe du rang entier $n$ à la fréquence $n/T$ (Hz si $T$ est en
    secondes) : seule l'échelle change, pas la forme du spectre. Le
    spectre bilatéral ne contient pas plus d'information que le spectre
    unilatéral -- $c_{-n}=\overline{c_n}$ pour un signal réel, donc $c_n$
    (pour $n\geq0$) suffit à tout reconstruire -- mais c'est la forme
    naturelle de l'écriture complexe, et celle qui généralisera
    directement à la transformée de Fourier (chapitre suivant).
\end{solution}

\q{Reconstruire le carré en ne gardant que les $c_n$ d'indice
$n\geq0$ (on retire tous les coefficients négatifs). Comparer au signal
original ainsi qu'à la reconstruction complète (symétrique). Que
devient la partie imaginaire ?}
\begin{solution}
    La reconstruction "asymétrique" n'est plus un signal réel : sa partie
    imaginaire n'est plus négligeable (elle est même du même ordre de
    grandeur que la partie réelle). En brisant la symétrie
    $c_{-n}=\overline{c_n}$, on perd la garantie que la somme reste
    réelle. C'est cette symétrie -- pas seulement la présence de tous les
    rangs -- qui assure que la série reconstruit bien un signal réel.
\end{solution}


%% ─── SECTION B ──────────────────────────────────────────────────
\section{Décalage temporel}

\subsection*{Objectif}
Retrouver, sur un exemple, la propriété de décalage temporel : que
devient le spectre d'un signal quand on le décale dans le temps ?

\q{Exécuter la cellule qui affiche $f(t)$ (le signal carré) puis
$g(t)=f(t+\varphi)$. Quel est le lien entre les deux tracés ?}
\begin{solution}
    $g$ est le tracé de $f$ translaté horizontalement de $\varphi$ (vers
    la gauche si $\varphi>0$) : même forme, même période, décalée dans le
    temps.
\end{solution}

\q{Quelles valeurs de $\varphi$ rend $f$ et $g$ identiques (sur toute la
droite réelle), et pourquoi ?}
\begin{solution}
    $\varphi = kT$, $k\in\mathbb Z$ (tout multiple de la période) :
    $f(t+kT)=f(t)$ pour tout $t$, exactement la définition de la
    $T$-périodicité de $f$.
\end{solution}

\q{\textbf{Sur papier.} Soit $f$ une fonction $T$-périodique de
coefficients complexes $(c_n)$. On pose $g(t)=f(t+\varphi)$. Exprimer les
coefficients complexes $(d_n)$ de $g$ en fonction de $(c_n)$.}
\begin{solution}
    $d_n = \dfrac{1}{T}\displaystyle\int_0^T f(t+\varphi)\,e^{-jn\omega_0 t}\,\mathrm dt$.
    Avec le changement de variable $u=t+\varphi$ (et $f$ $T$-périodique,
    donc l'intégrale sur une période ne dépend pas des bornes choisies) :
    $d_n=\dfrac1T\displaystyle\int f(u)\,e^{-jn\omega_0(u-\varphi)}\,\mathrm du
    = e^{jn\omega_0\varphi}\cdot\dfrac1T\displaystyle\int f(u)\,e^{-jn\omega_0 u}\,\mathrm du
    = c_n\, e^{jn\omega_0\varphi}$.
\end{solution}

\q{Exécuter la cellule qui calcule numériquement les coefficients de $f$
et de $g$, et les comparer. Vérifier l'accord avec la formule théorique trouver. %$d_n=c_n\,e^{jn\omega_0\varphi}$.
}
\begin{solution}
    Les valeurs numériques ($d_n$ calculés directement sur $g$) et
    prédites ($c_n\,e^{jn\omega_0\varphi}$) coïncident, aux erreurs
    d'arrondi près -- ce qui confirme la formule de B3.
\end{solution}

\q{Tracer les spectres bilatéraux de $f$ et de $g$. Que remarquez-vous ?}
\begin{solution}
    $|d_n|=|c_n\,e^{jn\omega_0\varphi}|=|c_n|$ pour tout $n$ (un facteur
    de module $1$ ne change pas le module) : les deux spectres
    \emph{d'amplitude} sont rigoureusement identiques. Un décalage
    temporel ne change que la \emph{phase} des coefficients, jamais leur
    module : c'est cohérent avec l'intuition physique -- décaler un
    signal dans le temps ne change ni sa forme ni son contenu fréquentiel,
    seulement l'instant où chaque motif apparaît.
\end{solution}


%% ─── SECTION C ──────────────────────────────────────────────────
\section{Superposition d'un signal et de son décalage}

\subsection*{Contexte}
On considère une tension $T$-périodique qui ne prend que $N$ valeurs
distinctes $v_0,\dots,v_{N-1}$, chacune maintenue pendant une durée
$T/N$ : un signal en escalier à pas régulier, sans lien de symétrie
particulier entre les paliers. On le note $f(t)$.

\begin{lesson}[Fonction \texttt{make\_levels}]
    \texttt{make\_levels(N, T, A)} construit $f$ : $N$ paliers de largeur
    $T/N$, de hauteurs tirées indépendamment dans $[-A,A]$.
\end{lesson}

\q{Afficher ce signal pour $N=4$. Ses coefficients complexes sont-ils, a
priori, réels, imaginaires purs, ou quelconques ? Pourquoi ?}
\begin{solution}
    Quelconques : rien n'impose de parité aux valeurs
    $v_0,\dots,v_{N-1}$, contrairement au signal carré (section A). Les
    $c_n$ n'ont donc, en général, ni partie réelle ni partie imaginaire
    nulle.
\end{solution}

\q{On pose $h(t)=f(t)+f(t+T/2)$. Afficher $h$. Quelle est sa période
apparente ?}
\begin{solution}
    $T/2$ : le motif de $h$ se répète deux fois sur $[0,T)$, visible sur
    le tracé.
\end{solution}

\q{\textbf{Sur papier.} Exprimer $c_n[h]$ en fonction de $c_n[f]$. Pour
quels rangs $n$ ce coefficient est-il nécessairement nul, quelle que
soit $f$ ? Est-ce cohérent avec la période observée en C2 ?}
\begin{solution}
    Par linéarité et décalage (B3), avec $\omega_0(T/2)=\pi$ :
    $$
    c_n[h] = c_n[f]\left(1+e^{jn\pi}\right) = c_n[f]\left(1+(-1)^n\right),
    $$
    soit $2c_n[f]$ si $n$ est pair, $0$ si $n$ est impair -- quelle que
    soit $f$. Cohérent avec C2 : un signal dont seuls les rangs pairs
    (relatifs à $T$) sont non nuls s'écrit avec des exponentielles de
    fréquences multiples de $2\omega_0$, donc de période (au plus) $T/2$.
\end{solution}

\q{Calculer numériquement $c_n[f]$ et $c_n[h]$. Les rangs impairs
s'annulent-ils effectivement ?}
\begin{solution}
    Oui, à la précision numérique près, pour tout rang impair calculé --
    quelles que soient les valeurs $v_0,\dots,v_{N-1}$ tirées.
\end{solution}

\q{Tracer les spectres bilatéraux de $f$ et de $h$. Conclure.}
\begin{solution}
    Le spectre de $h$ ne conserve que les raies paires de celui de $f$
    (doublées en amplitude) ; les raies impaires ont disparu -- la
    signature spectrale exacte d'un signal deux fois plus rapide.
\end{solution}

\q{Recommencer l'opération sur $h$, avec un décalage égal à la moitié de
sa période apparente ($T/4$). Calculer $h_2(t)=h(t)+h(t+T/4)$, tracer son
spectre et son signal temporel. Que devient $h_2$ ?}
\begin{solution}
    $h_2$ ne conserve que le rang $n=0$ (tous les autres rangs calculés
    sont numériquement nuls) : c'est un signal constant, égal à $4$ fois
    la moyenne de $f$.
\end{solution}

\q{Plus généralement, que suggère cette construction sur ce qu'on peut
obtenir d'un signal en l'additionnant seulement à des versions décalées
de lui-même ?}
\begin{solution}
    En choisissant soigneusement des décalages successifs, la simple
    addition de copies retardées d'un signal permet d'effacer ou de
    conserver des portions choisies de son spectre -- jusqu'à l'annuler
    presque entièrement, comme ici. C'est un principe de construction
    (combiner des copies décalées pour obtenir un effet voulu sur le
    spectre), pas seulement un résultat de calcul.
\end{solution}

\q{\textbf{Bonus (facultatif).} Démontrer le résultat de C6 : exprimer
$c_n[h_2]$ en fonction de $c_n[f]$, et montrer que seuls les rangs
multiples de $4$ peuvent être non nuls. Pourquoi ce résultat entraîne-t-il,
pour ce signal $f$ à $N=4$ paliers, que $h_2$ est constant ?}
\begin{solution}
    En combinant B3 deux fois -- une fois avec $\varphi=T/2$ (C3), une
    fois avec $T/4$ --~:
    $$
    c_n[h_2] = c_n[f]\left(1+(-1)^n\right)\left(1+e^{jn\pi/2}\right).
    $$
    Le premier facteur impose $n$ pair ; le second s'annule pour
    $e^{jn\pi/2}=-1$, soit $n\equiv2\pmod4$. Seuls les rangs multiples de
    $4$ peuvent donc survivre. Or, pour un signal à $N=4$ paliers de
    largeur égale, $c_n[f]=0$ automatiquement pour tout multiple de $4$
    (chaque palier contribue une intégrale de l'exponentielle sur un
    cycle complet, qui s'annule) : il ne reste donc que $n=0$, la moyenne
    de $f$. Pour un nombre de paliers qui n'est pas une puissance de $2$,
    cette coïncidence n'a pas lieu et $h_2$ n'est en général pas
    constante.
\end{solution}

\end{document}
